\input{preamble}

% OK, start here.
%
\begin{document}

\title{Espaces algébriques décents}


\maketitle

\phantomsection
\label{section-phantom}

\tableofcontents

\section{Introduction}
\label{section-introduction}

\noindent
Dans ce chapitre, nous étudions les propriétés « locales » des espaces
algébriques généraux, c'est-à-dire de ceux qui ne sont pas quasi-séparés.
Les espaces algébriques quasi-séparés sont étudiés dans \cite{Kn}.
Il apparaît que des phénomènes essentiellement nouveaux se produisent,
notamment en ce qui concerne les points et leurs spécialisations, pour les
espaces algébriques plus généraux. D'autre part, dans la plupart des résultats
fondamentaux sur les espaces algébriques, il n'est pas nécessaire de se
préoccuper de ces phénomènes. C'est pourquoi nous avons choisi de réunir cette
matière dans un chapitre séparé, après l'exposé habituel de la théorie.



\section{Conventions}
\label{section-conventions}

\noindent
Nous supposons une fois pour toutes que tous les schémas sont contenus dans
un grand site fppf $\Sch_{fppf}$. De plus, tout anneau $A$ considéré est tel
que $\Spec(A)$ est (isomorphe à) un objet de ce grand site.

\medskip\noindent
Soit $S$ un schéma et soit $X$ un espace algébrique sur $S$.
Dans ce chapitre et le suivant, nous noterons $X \times_S X$ le produit de $X$
par lui-même (dans la catégorie des espaces algébriques sur $S$), au lieu de
$X \times X$.



\section{Fibres universellement bornées}
\label{section-universally-bounded}

\noindent
Nous examinons brièvement ce que signifie, pour un morphisme d'un schéma vers
un espace algébrique, le fait d'avoir des fibres universellement bornées. Voir
Morphisms, Section \ref{morphisms-section-universally-bounded}
pour les définitions et résultats analogues concernant les morphismes de
schémas.

\begin{definition}
\label{definition-universally-bounded}
Soit $S$ un schéma. Soit $X$ un espace algébrique sur $S$, et soit
$U$ un schéma sur $S$. Soit $f : U \to X$ un morphisme sur $S$.
Nous dirons que {\it les fibres de $f$ sont universellement
bornées}\footnote{Cette terminologie n'est probablement pas standard.}
s'il existe un entier $n$ tel que, pour tout corps $k$ et tout morphisme
$\Spec(k) \to X$, le produit fibré $\Spec(k) \times_X U$ soit un schéma fini
sur $k$, dont le degré sur $k$ est $\leq n$.
\end{definition}

\noindent
Cette définition a un sens, car le produit fibré
$\Spec(k) \times_X U$ est un schéma. De plus, si $X$ est un schéma,
on retrouve la notion de
Morphisms, Definition \ref{morphisms-definition-universally-bounded}
en vertu de
Morphisms, Lemma \ref{morphisms-lemma-characterize-universally-bounded}.

\begin{lemma}
\label{lemma-composition-universally-bounded}
Soit $S$ un schéma. Soit $X$ un espace algébrique sur $S$.
Soit $V \to U$ un morphisme de schémas sur $S$, et soit
$U \to X$ un morphisme de $U$ vers $X$. Si les fibres de
$V \to U$ et de $U \to X$ sont universellement bornées, alors celles de
$V \to X$ le sont aussi.
\end{lemma}

\begin{proof}
Soit $n$ un entier qui convient à $V \to U$, et soit $m$ un entier qui
convient à $U \to X$ dans la
Définition \ref{definition-universally-bounded}.
Soit $\Spec(k) \to X$ un morphisme, où $k$ est un corps.
Considérons les morphismes
$$
\Spec(k) \times_X V
\longrightarrow
\Spec(k) \times_X U
\longrightarrow
\Spec(k).
$$
Par hypothèse, le schéma $\Spec(k) \times_X U$ est fini de degré au plus $m$
sur $k$, et $n$ borne le degré des fibres du premier morphisme. Par
Morphisms, Lemma \ref{morphisms-lemma-composition-universally-bounded},
nous en déduisons que $\Spec(k) \times_X V$ est fini sur $k$ de degré au plus
$nm$.
\end{proof}

\begin{lemma}
\label{lemma-base-change-universally-bounded}
Soit $S$ un schéma.
Soit $Y \to X$ un morphisme représentable d'espaces algébriques sur $S$.
Soit $U \to X$ un morphisme d'un schéma vers $X$.
Si les fibres de $U \to X$ sont universellement bornées, alors celles de
$U \times_X Y \to Y$ le sont aussi.
\end{lemma}

\begin{proof}
Cela résulte immédiatement de la définition et des propriétés des produits
fibrés. (Notons que $U \times_X Y$ est un schéma puisque nous avons supposé
$Y \to X$ représentable, si bien que la définition s'applique.)
\end{proof}

\begin{lemma}
\label{lemma-descent-universally-bounded}
Soit $S$ un schéma. Soit $g : Y \to X$ un morphisme représentable d'espaces
algébriques sur $S$. Soit $f : U \to X$ un morphisme d'un schéma vers $X$.
Soit $f' : U \times_X Y \to Y$ le changement de base de $f$.
Si
$$
\Im(|f| : |U| \to |X|) \subset \Im(|g| : |Y| \to |X|)
$$
et si $f'$ a des fibres universellement bornées, alors $f$ a des fibres
universellement bornées.
\end{lemma}

\begin{proof}
Soit $n \geq 0$ un entier qui borne les degrés des produits fibrés
$\Spec(k) \times_Y (U \times_X Y)$ au sens de la
Définition \ref{definition-universally-bounded} appliquée au morphisme $f'$.
Nous affirmons que $n$ convient aussi à $f$. En effet, supposons que
$x : \Spec(k) \to X$ soit un morphisme depuis le spectre d'un corps.
Alors, soit $\Spec(k) \times_X U$ est vide (et il n'y a rien à démontrer),
soit $x$ appartient à l'image de $|f|$. D'après
Properties of Spaces,
Lemma \ref{spaces-properties-lemma-points-cartesian}
et l'hypothèse du lemme, cela signifie qu'il existe une extension de corps
$k'/k$ et un diagramme commutatif
$$
\xymatrix{
\Spec(k') \ar[r] \ar[d] & Y \ar[d] \\
\Spec(k) \ar[r] & X
}
$$
On obtient alors
$$
\Spec(k') \times_Y (U \times_X Y) =
\Spec(k') \times_{\Spec(k)} (\Spec(k) \times_X U)
$$
Comme le schéma $\Spec(k') \times_Y (U \times_X Y)$ est supposé fini de degré
$\leq n$ sur $k'$, il en résulte que $\Spec(k) \times_X U$ est lui aussi fini
de degré $\leq n$ sur $k$, comme voulu. (Nous omettons quelques détails.)
\end{proof}

\begin{lemma}
\label{lemma-universally-bounded-permanence}
Soit $S$ un schéma. Soit $X$ un espace algébrique sur $S$.
Considérons un diagramme commutatif
$$
\xymatrix{
U \ar[rd]_g \ar[rr]_f & & V \ar[ld]^h \\
& X &
}
$$
où $U$ et $V$ sont des schémas. Si $g$ a des fibres universellement bornées
et si $f$ est surjectif et plat, alors $h$ a lui aussi des fibres
universellement bornées.
\end{lemma}

\begin{proof}
Supposons que $g$ ait des fibres universellement bornées et que $f$ soit
surjectif et plat. Soit $n \geq 0$ un entier qui borne les degrés des schémas
$\Spec(k) \times_X U$ au sens de la
Définition \ref{definition-universally-bounded}.
Nous affirmons que $n$ convient aussi à $h$.
Soit $\Spec(k) \to X$ un morphisme du spectre d'un corps vers $X$.
Considérons le morphisme de schémas
$$
\Spec(k) \times_X U \longrightarrow \Spec(k) \times_X V
$$
Il est plat et surjectif. Par hypothèse, le schéma de gauche est fini de degré
$\leq n$ sur $\Spec(k)$. Il résulte de
Morphisms, Lemma \ref{morphisms-lemma-universally-bounded-permanence}
que le degré du schéma de droite est lui aussi majoré par $n$, comme voulu.
\end{proof}

\begin{lemma}
\label{lemma-universally-bounded-finite-fibres}
Soit $S$ un schéma.
Soit $X$ un espace algébrique sur $S$, et soit $U$ un schéma sur $S$.
Soit $\varphi : U \to X$ un morphisme sur $S$.
Si les fibres de $\varphi$ sont universellement bornées, il existe un entier
$n$ tel que chaque fibre de $|U| \to |X|$ compte au plus $n$ éléments.
\end{lemma}

\begin{proof}
L'entier $n$ de la Définition \ref{definition-universally-bounded} convient.
En effet, choisissons $x \in |X|$. Représentons $x$ par un morphisme
$x : \Spec(k) \to X$. Nous obtenons alors un diagramme commutatif
$$
\xymatrix{
\Spec(k) \times_X U \ar[r] \ar[d] & U \ar[d] \\
\Spec(k) \ar[r]^x & X
}
$$
qui montre (au moyen de
Properties of Spaces,
Lemma \ref{spaces-properties-lemma-points-cartesian})
que l'image réciproque de $x$ dans $|U|$ est l'image de la flèche horizontale
supérieure. Puisque $\Spec(k) \times_X U$ est fini de degré $\leq n$ sur $k$,
il possède au plus $n$ points.
\end{proof}








\section{Conditions de finitude et points}
\label{section-points-monomorphisms}

\noindent
Dans cette section, nous approfondissons la question de savoir quand des
points peuvent être représentés par des monomorphismes de spectres de corps
dans l'espace.

\begin{remark}
\label{remark-recall}
Avant de démontrer le lemme suivant, rappelons quelques faits sur les
morphismes étales de schémas :
\begin{enumerate}
\item Un morphisme étale est plat; les générisations se relèvent donc le long
d'un morphisme étale
(Morphisms, Lemmas \ref{morphisms-lemma-etale-flat}
et \ref{morphisms-lemma-generalizations-lift-flat}).
\item Un morphisme étale est non ramifié; un morphisme non ramifié est
localement quasi-fini; ses fibres sont donc discrètes
(Morphisms, Lemmas \ref{morphisms-lemma-flat-unramified-etale},
\ref{morphisms-lemma-unramified-quasi-finite}, et
\ref{morphisms-lemma-quasi-finite-at-point-characterize}).
\item Un morphisme étale quasi-compact est quasi-fini et possède en particulier
des fibres finies
(Morphisms, Lemmas \ref{morphisms-lemma-quasi-finite-locally-quasi-compact} et
\ref{morphisms-lemma-quasi-finite}).
\item Un schéma étale sur un corps $k$ est une réunion disjointe de spectres
d'extensions finies séparables de $k$
(Morphisms, Lemma \ref{morphisms-lemma-etale-over-field}).
\end{enumerate}
Pour une étude générale des morphismes étales, voir
\'Etale Morphisms, Section \ref{etale-section-etale-morphisms}.
\end{remark}

\begin{lemma}
\label{lemma-U-finite-above-x}
Soit $S$ un schéma. Soit $X$ un espace algébrique sur $S$.
Soit $x \in |X|$. Les conditions suivantes sont équivalentes :
\begin{enumerate}
\item il existe une famille de schémas $U_i$ et de morphismes étales
$\varphi_i : U_i \to X$ telle que
$\coprod \varphi_i : \coprod U_i \to X$ soit surjectif,
et telle que, pour tout $i$, la fibre de
$|U_i| \to |X|$ au-dessus de $x$ soit finie, et
\item pour tout schéma affine $U$ et tout morphisme étale $\varphi : U \to X$,
la fibre de $|U| \to |X|$ au-dessus de $x$ est finie.
\end{enumerate}
\end{lemma}

\begin{proof}
L'implication (2) $\Rightarrow$ (1) est immédiate.
Soit $\varphi_i : U_i \to X$ une famille de morphismes étales comme en (1).
Soit $\varphi : U \to X$ un morphisme étale d'un schéma affine vers $X$.
Considérons les diagrammes de produits fibrés
$$
\xymatrix{
U \times_X U_i \ar[r]_-{p_i} \ar[d]_{q_i} & U_i \ar[d]^{\varphi_i} \\
U \ar[r]^\varphi & X
}
\quad \quad
\xymatrix{
\coprod U \times_X U_i \ar[r]_-{\coprod p_i} \ar[d]_{\coprod q_i} &
\coprod U_i \ar[d]^{\coprod \varphi_i} \\
U \ar[r]^\varphi & X
}
$$
Comme $q_i$ est étale, il est ouvert (voir la Remarque \ref{remark-recall}).
De plus, le morphisme $\coprod q_i$ est surjectif.
Il existe donc un nombre fini d'indices $i_1, \ldots, i_n$ et des ouverts
quasi-compacts $W_{i_j} \subset U \times_X U_{i_j}$ dont les images recouvrent
$U$.
Le morphisme $p_i$ est étale, donc localement quasi-fini (voir la remarque sur
les morphismes étales ci-dessus). Nous pouvons donc appliquer
Morphisms, Lemma
\ref{morphisms-lemma-locally-quasi-finite-qc-source-universally-bounded}
pour voir que les fibres de
$p_{i_j}|_{W_{i_j}} : W_{i_j} \to U_{i_j}$ sont finies.
D'après
Properties of Spaces,
Lemma \ref{spaces-properties-lemma-points-cartesian}
et l'hypothèse sur $\varphi_i$, nous en concluons que la fibre de $\varphi$
au-dessus de $x$ est finie. Autrement dit, (2) est vérifiée.
\end{proof}

\begin{lemma}
\label{lemma-R-finite-above-x}
Soit $S$ un schéma. Soit $X$ un espace algébrique sur $S$.
Soit $x \in |X|$. Les conditions suivantes sont équivalentes :
\begin{enumerate}
\item il existe un schéma $U$, un morphisme étale
$\varphi : U \to X$, et des points $u, u' \in U$ dont l'image est
$x$ tels que, si l'on pose $R = U \times_X U$, la fibre de
$$
|R| \to |U| \times_{|X|} |U|
$$
au-dessus de $(u, u')$ soit finie,
\item pour tout schéma $U$, tout morphisme étale $\varphi : U \to X$ et
tous points $u, u' \in U$ dont l'image est
$x$, si l'on pose $R = U \times_X U$, la fibre de
$$
|R| \to |U| \times_{|X|} |U|
$$
au-dessus de $(u, u')$ est finie,
\item il existe un morphisme $\Spec(k) \to X$, où $k$ est un corps,
dans la classe d'équivalence de $x$, tel que les projections
$\Spec(k) \times_X \Spec(k) \to \Spec(k)$ soient
étales et quasi-compactes, et
\item il existe un monomorphisme $\Spec(k) \to X$, où $k$ est un corps,
dans la classe d'équivalence de $x$.
\end{enumerate}
\end{lemma}

\begin{proof}
Supposons (1), c'est-à-dire soit $\varphi : U \to X$ un morphisme étale
d'un schéma vers $X$, et soient $u, u'$ des points de $U$ situés au-dessus de
$x$ tels que la fibre de $|R| \to |U| \times_{|X|} |U|$ au-dessus de $(u, u')$
soit un ensemble fini. Dans cette démonstration, nous considérons un point
$u = \Spec(\kappa(u))$ comme un schéma. Notons que $u \to U$, $u' \to U$ sont
des monomorphismes (voir
Schemes, Lemma \ref{schemes-lemma-injective-points-surjective-stalks}); ainsi,
$u \times_X u' \to R = U \times_X U$ est un monomorphisme.
Dans ce langage, l'hypothèse signifie exactement que
$u \times_X u'$ est un schéma dont l'espace topologique sous-jacent possède
un nombre fini de points.
Soit $\psi : W \to X$ un morphisme étale d'un schéma vers $X$.
Soient $w, w' \in W$ des points de $W$ dont l'image est $x$.
Nous devons montrer que $w \times_X w'$ est un schéma dont l'espace topologique
sous-jacent possède un nombre fini de points.
Considérons le diagramme de produit fibré
$$
\xymatrix{
W \times_X U \ar[r]_p \ar[d]_q & U \ar[d]^\varphi \\
W \ar[r]^\psi & X
}
$$
Comme $x$ est l'image de $u$ et de $u'$, nous pouvons choisir des points
$\tilde w, \tilde w'$ de $W \times_X U$ tels que $q(\tilde w) = w$,
$q(\tilde w') = w'$, $u = p(\tilde w)$ et $u' = p(\tilde w')$; voir
Properties of Spaces,
Lemma \ref{spaces-properties-lemma-points-cartesian}.
Comme $p$ et $q$ sont étales, les extensions de corps
$\kappa(w) \subset \kappa(\tilde w) \supset \kappa(u)$ et
$\kappa(w') \subset \kappa(\tilde w') \supset \kappa(u')$ sont
finies séparables; voir la Remarque \ref{remark-recall}.
Nous obtenons alors un diagramme commutatif
$$
\xymatrix{
w \times_X w' \ar[d] &
\tilde w \times_X \tilde w' \ar[l] \ar[d] \ar[r] &
u \times_X u' \ar[d] \\
w \times_X w' &
\tilde w \times_S \tilde w' \ar[l] \ar[r] &
u \times_S u'
}
$$
où les carrés sont des carrés de produits fibrés. Les morphismes
horizontaux inférieurs sont étales et quasi-compacts, car tout schéma de la
forme $\Spec(k) \times_S \Spec(k')$ est affine, et d'après nos observations
ci-dessus sur les extensions de corps.
Nous voyons donc que les flèches horizontales supérieures sont étales et
quasi-compactes, et qu'elles ont par conséquent des fibres finies.
Nous avons vu plus haut que $|u \times_X u'|$ est fini; nous en concluons que
$|w \times_X w'|$ est fini. Autrement dit, (2) est vérifiée.

\medskip\noindent
Supposons (2). Soit $U \to X$ un morphisme étale d'un schéma $U$
tel que $x$ appartienne à l'image de $|U| \to |X|$. Soit $u \in U$ un
point dont l'image est $x$. Nous avons alors vu au paragraphe précédent que
$u = \Spec(\kappa(u)) \to X$ a la propriété que
$u \times_X u$ possède un espace topologique sous-jacent fini. D'autre part,
les projections $u \times_X u \to u$ sont les composés
$$
u \times_X u \longrightarrow
u \times_X U \longrightarrow
u \times_X X = u,
$$
c'est-à-dire les composés d'un monomorphisme (le changement de base du
monomorphisme $u \to U$) et d'un morphisme étale (le changement de base du
morphisme étale $U \to X$). Ainsi, $u \times_X U$ est une réunion disjointe
de spectres de corps qui sont des extensions finies séparables de $\kappa(u)$
(voir la Remarque \ref{remark-recall}). Puisque $u \times_X u$ est fini, son
image dans $u \times_X U$ est une réunion disjointe finie de spectres de corps
qui sont des extensions finies séparables de $\kappa(u)$. Par
Schemes, Lemma \ref{schemes-lemma-mono-towards-spec-field},
nous en concluons que $u \times_X u$ est une réunion disjointe finie de
spectres de corps qui sont des extensions finies séparables de $\kappa(u)$.
Autrement dit, nous voyons que $u \times_X u \to u$ est quasi-compact et
étale. Cela signifie que (3) est vérifiée.

\medskip\noindent
Démontrons que (3) implique (4). Soit $\Spec(k) \to X$ un morphisme
du spectre d'un corps vers $X$, dans la classe d'équivalence de $x$,
tel que les deux projections
$t, s : R = \Spec(k) \times_X \Spec(k)  \to \Spec(k)$
soient quasi-compactes et étales.
Cela signifie en particulier que $R$ est une relation d'équivalence étale
sur $\Spec(k)$.
Par Spaces, Theorem \ref{spaces-theorem-presentation},
nous savons que le faisceau quotient
$X' = \Spec(k)/R$ est un espace algébrique. Par
Groupoids, Lemma \ref{groupoids-lemma-quotient-groupoid-restrict},
le morphisme $X' \to X$ est un monomorphisme.
Comme $s, t$ sont quasi-compacts, nous voyons que $R$ est quasi-compact;
Properties of Spaces,
Lemma \ref{spaces-properties-lemma-point-like-spaces}
s'applique donc à $X'$, et nous voyons que
$X' = \Spec(k')$ pour un certain corps $k'$. Nous obtenons ainsi une
factorisation
$$
\Spec(k) \longrightarrow
\Spec(k') \longrightarrow X
$$
qui montre que $\Spec(k') \to X$ est un monomorphisme dont l'image est
$x \in |X|$. Autrement dit, (4) est vérifiée.

\medskip\noindent
Enfin, démontrons que (4) implique (1). Soit $\Spec(k) \to X$
un monomorphisme, où $k$ est un corps, dans la classe d'équivalence de $x$.
Soit $U \to X$ un morphisme étale surjectif d'un schéma $U$ vers $X$.
Soit $u \in U$ un point au-dessus de $x$. Puisque $\Spec(k) \times_X u$
est non vide et que $\Spec(k) \times_X u \to u$ est un monomorphisme,
nous en concluons que $\Spec(k) \times_X u = u$ (voir
Schemes, Lemma \ref{schemes-lemma-mono-towards-spec-field}).
Le morphisme $u \to U \to X$ se factorise donc par $\Spec(k) \to X$; voici
le diagramme
$$
\xymatrix{
u \ar[r] \ar[d] & U \ar[d] \\
\Spec(k) \ar[r] & X
}
$$
Comme la flèche verticale de droite est étale, il en résulte que
$\kappa(u)/k$ est une extension finie séparable. Par conséquent,
$$
u \times_X u = u \times_{\Spec(k)} u
$$
est un schéma fini, et le résultat découle de l'analyse de la propriété (1)
faite au premier paragraphe de cette démonstration.
\end{proof}

\begin{lemma}
\label{lemma-weak-UR-finite-above-x}
Soit $S$ un schéma. Soit $X$ un espace algébrique sur $S$.
Soit $x \in |X|$.
Soit $U$ un schéma et soit $\varphi : U \to X$ un morphisme étale.
Les conditions suivantes sont équivalentes :
\begin{enumerate}
\item $x$ appartient à l'image de $|U| \to |X|$ et, si l'on pose
$R = U \times_X U$, les fibres des deux applications
$$
|U| \longrightarrow |X|
\quad\text{and}\quad
|R| \longrightarrow |X|
$$
au-dessus de $x$ sont finies,
\item il existe un monomorphisme $\Spec(k) \to X$, où $k$ est un corps,
dans la classe d'équivalence de $x$, et
le produit fibré $\Spec(k) \times_X U$ est
un schéma fini non vide sur $k$.
\end{enumerate}
\end{lemma}

\begin{proof}
Supposons (1). Cette condition implique clairement la première condition du
Lemme \ref{lemma-R-finite-above-x}; nous obtenons donc un monomorphisme
$\Spec(k) \to X$ dans la classe de $x$. En prenant le produit fibré,
nous voyons que $\Spec(k) \times_X U \to \Spec(k)$ est un schéma
étale sur $\Spec(k)$ ayant un nombre fini de points, donc un schéma fini
non vide sur $k$, c'est-à-dire que (2) est vérifiée.

\medskip\noindent
Supposons (2). Par hypothèse, $x$ appartient à l'image de
$|U| \to |X|$. La finitude de la fibre de
$|U| \to |X|$ au-dessus de $x$ est claire, puisque cette fibre est égale à
$|\Spec(k) \times_X U|$ d'après
Properties of Spaces,
Lemma \ref{spaces-properties-lemma-points-cartesian}.
La finitude de la fibre de $|R| \to |X|$ au-dessus de $x$ est également
claire, puisqu'elle est égale à l'ensemble sous-jacent au schéma
$$
(\Spec(k) \times_X U) \times_{\Spec(k)} (\Spec(k) \times_X U)
$$
qui est fini sur $k$. Ainsi, (1) est vérifiée.
\end{proof}

\begin{lemma}
\label{lemma-UR-finite-above-x}
Soit $S$ un schéma. Soit $X$ un espace algébrique sur $S$.
Soit $x \in |X|$. Les conditions suivantes sont équivalentes :
\begin{enumerate}
\item pour tout schéma affine $U$ et tout morphisme étale
$\varphi : U \to X$, si l'on pose $R = U \times_X U$, les fibres des deux
applications
$$
|U| \longrightarrow |X|
\quad\text{and}\quad
|R| \longrightarrow |X|
$$
au-dessus de $x$ sont finies,
\item il existe des schémas $U_i$ et des morphismes étales
$U_i \to X$ tels que $\coprod U_i \to X$ soit surjectif et que, pour tout
$i$, si l'on pose $R_i = U_i \times_X U_i$, les fibres des deux applications
$$
|U_i| \longrightarrow |X|
\quad\text{and}\quad
|R_i| \longrightarrow |X|
$$
au-dessus de $x$ soient finies,
\item il existe un monomorphisme $\Spec(k) \to X$, où $k$ est un corps,
dans la classe d'équivalence de $x$, et, pour tout schéma affine $U$ et tout
morphisme étale $U \to X$, le produit fibré $\Spec(k) \times_X U$ est
un schéma fini sur $k$,
\item il existe un monomorphisme quasi-compact $\Spec(k) \to X$,
où $k$ est un corps dans la classe d'équivalence de $x$,
\item il existe un morphisme quasi-compact $\Spec(k) \to X$,
où $k$ est un corps dans la classe d'équivalence de $x$, et
\item tout morphisme $\Spec(k) \to X$, où $k$ est un corps dans la
classe d'équivalence de $x$, est quasi-compact.
\end{enumerate}
\end{lemma}

\begin{proof}
L'équivalence de (1) et (3) résulte de l'application du
Lemme \ref{lemma-weak-UR-finite-above-x}
à tout morphisme étale $U \to X$ avec $U$ affine.
Il est clair que (3) implique (2).
Supposons que $U_i \to X$ et $R_i$ soient comme en (2). Le
Lemme \ref{lemma-U-finite-above-x}
nous permet de conclure que, pour tout schéma affine $U$ et tout morphisme
étale $U \to X$, la fibre de $|U| \to |X|$ au-dessus de $x$ est finie.
Écrivons cette fibre $\{u_1, \ldots, u_n\}$. Comme la condition (1) du
Lemme \ref{lemma-R-finite-above-x}
s'applique à $U_i \to X$ pour un certain $i$ tel que $x$ appartienne à
l'image de $|U_i| \to |X|$, nous voyons que la fibre de
$|R = U \times_X U| \to |U| \times_{|X|} |U|$
est finie au-dessus de $(u_a, u_b)$, où $a, b \in \{1, \ldots, n\}$.
La fibre de $|R| \to |X|$ au-dessus de $x$ est donc finie.
Nous voyons ainsi que (1) est vérifiée. À ce stade, nous savons que
(1), (2) et (3) sont équivalentes.

\medskip\noindent
Si (4) est vérifiée, alors, pour tout schéma affine $U$ et tout morphisme
étale $U \to X$, le schéma $\Spec(k) \times_X U$ est d'une part
étale sur $k$ (donc une réunion disjointe de spectres d'extensions finies
séparables de $k$, d'après la
Remarque \ref{remark-recall})
et d'autre part quasi-compact sur $U$ (donc quasi-compact).
Nous voyons ainsi que (3) est vérifiée.
Réciproquement, si $U_i \to X$ est comme en (2) et si $\Spec(k) \to X$
est un monomorphisme comme en (3), alors
$$
\coprod \Spec(k) \times_X U_i
\longrightarrow
\coprod U_i
$$
est quasi-compact (car, au-dessus de chaque $U_i$, nous voyons que
$\Spec(k) \times_X U_i$ est une réunion disjointe finie de spectres de corps).
Ainsi, $\Spec(k) \to X$ est quasi-compact d'après
Morphisms of Spaces, Lemma \ref{spaces-morphisms-lemma-quasi-compact-local}.

\medskip\noindent
Il est immédiat que (4) implique (5). Réciproquement, soit
$\Spec(k) \to X$ un morphisme quasi-compact dans la classe d'équivalence de
$x$. Soit $U \to X$ un morphisme étale avec $U$ affine. Considérons le
produit fibré
$$
\xymatrix{
F \ar[r] \ar[d] & U \ar[d] \\
\Spec(k) \ar[r] & X
}
$$
Alors $F \to U$ est quasi-compact, donc $F$ est quasi-compact.
D'autre part, $F \to \Spec(k)$ est étale; ainsi, $F$ est une
réunion disjointe finie de spectres d'extensions finies séparables de $k$
(Remarque \ref{remark-recall}). Comme l'image de $|F| \to |U|$
est la fibre de $|U| \to |X|$ au-dessus de $x$ (Properties of Spaces, Lemma
\ref{spaces-properties-lemma-points-cartesian}), nous en concluons que
la fibre de $|U| \to |X|$ au-dessus de $x$ est finie. Le schéma
$F \times_{\Spec(k)} F$ est aussi une réunion finie de spectres de corps,
car il est lui aussi quasi-compact et étale sur $\Spec(k)$.
Il existe un monomorphisme
$F \times_X F \to F \times_{\Spec(k)} F$; par conséquent, $F \times_X F$ est
une réunion disjointe finie de spectres de corps
(Schemes, Lemma \ref{schemes-lemma-mono-towards-spec-field}).
L'image de $F \times_X F \to U \times_X U = R$ est donc finie.
Comme cette image est la fibre de $|R| \to |X|$ au-dessus de $x$ d'après
Properties of Spaces, Lemma \ref{spaces-properties-lemma-points-cartesian},
nous en concluons que (1) est vérifiée. À ce stade, nous savons que
(1) -- (5) sont équivalentes.

\medskip\noindent
Il est clair que (6) implique (5). Réciproquement, supposons que
$\Spec(k) \to X$ soit comme en (4) et soit $\Spec(k') \to X$ un autre
morphisme, où $k'$ est un corps dans la classe d'équivalence de $x$. Par
Properties of Spaces, Lemma
\ref{spaces-properties-lemma-equivalence-class-point-monomorphism},
nous avons une factorisation $\Spec(k') \to \Spec(k) \to X$ du morphisme
donné. C'est un composé de morphismes quasi-compacts, donc un morphisme
quasi-compact (Morphisms of Spaces,
Lemma \ref{spaces-morphisms-lemma-composition-quasi-compact}), comme voulu.
\end{proof}

\begin{lemma}
\label{lemma-U-universally-bounded}
Soit $S$ un schéma. Soit $X$ un espace algébrique sur $S$.
Les conditions suivantes sont équivalentes :
\begin{enumerate}
\item il existe des schémas $U_i$ et des morphismes étales
$U_i \to X$ tels que $\coprod U_i \to X$ soit surjectif et que
chaque $U_i \to X$ ait des fibres universellement bornées, et
\item pour tout schéma affine $U$ et tout morphisme étale $\varphi : U \to X$,
les fibres de $U \to X$ sont universellement bornées.
\end{enumerate}
\end{lemma}

\begin{proof}
L'implication (2) $\Rightarrow$ (1) est immédiate.
Supposons (1). Soit $(\varphi_i : U_i \to X)_{i \in I}$ une famille de
morphismes étales de schémas vers $X$, qui recouvre $X$, telle que
chaque $\varphi_i$ ait des fibres universellement bornées.
Soit $\psi : U \to X$ un morphisme étale d'un schéma affine vers $X$.
Pour tout $i$, considérons le diagramme de produit fibré
$$
\xymatrix{
U \times_X U_i \ar[r]_{p_i} \ar[d]_{q_i} & U_i \ar[d]^{\varphi_i} \\
U \ar[r]^\psi & X
}
$$
Comme $q_i$ est étale, il est ouvert (voir la Remarque \ref{remark-recall}).
De plus, nous avons $U = \bigcup \Im(q_i)$, puisque la famille
$(\varphi_i)_{i \in I}$ est surjective. Comme $U$ est affine, donc
quasi-compact, nous pouvons trouver un nombre fini d'indices
$i_1, \ldots, i_n \in I$ et des ouverts quasi-compacts
$W_j \subset U \times_X U_{i_j}$ tels que
$U = \bigcup q_{i_j}(W_j)$.
Le morphisme $p_{i_j}$ est étale, donc localement quasi-fini
(voir la remarque sur les morphismes étales ci-dessus). Nous pouvons donc
appliquer Morphisms, Lemma
\ref{morphisms-lemma-locally-quasi-finite-qc-source-universally-bounded}
pour voir que les fibres de $p_{i_j}|_{W_j} : W_j \to U_{i_j}$ sont
universellement bornées. Par le
Lemme \ref{lemma-composition-universally-bounded},
nous voyons donc que les fibres de $W_j \to X$ sont universellement bornées.
Ainsi, $\coprod_{j = 1, \ldots, n} W_j \to X$ a lui aussi des fibres
universellement bornées. Comme $\coprod_{j = 1, \ldots, n} W_j \to X$ se
factorise par le morphisme étale surjectif
$\coprod q_{i_j}|_{W_j} : \coprod_{j = 1, \ldots, n} W_j \to U$, nous
voyons que les fibres de $U \to X$ sont universellement bornées d'après le
Lemme \ref{lemma-universally-bounded-permanence}.
Autrement dit, (2) est vérifiée.
\end{proof}

\begin{lemma}
\label{lemma-characterize-very-reasonable}
Soit $S$ un schéma.
Soit $X$ un espace algébrique sur $S$.
Les conditions suivantes sont équivalentes :
\begin{enumerate}
\item il existe un recouvrement de Zariski $X = \bigcup X_i$ et, pour
chaque $i$, un schéma $U_i$ ainsi qu'un morphisme étale surjectif
quasi-compact $U_i \to X_i$, et
\item il existe des schémas $U_i$ et des morphismes étales $U_i \to X$
tels que les projections $U_i \times_X U_i \to U_i$ soient quasi-compactes
et que $\coprod U_i \to X$ soit surjectif.
\end{enumerate}
\end{lemma}

\begin{proof}
Si (1) est vérifiée, alors les morphismes $U_i \to X_i \to X$ sont étales
(combiner Morphisms, Lemma \ref{morphisms-lemma-composition-etale}
et
Spaces, Lemmas
\ref{spaces-lemma-composition-representable-transformations-property} et
\ref{spaces-lemma-morphism-schemes-gives-representable-transformation-property}
).
De plus, puisque $U_i \times_X U_i = U_i \times_{X_i} U_i$,
les deux projections $U_i \times_X U_i \to U_i$ sont quasi-compactes.

\medskip\noindent
Si (2) est vérifiée, soit $X_i \subset X$ le sous-espace ouvert correspondant
à l'image de l'application ouverte $|U_i| \to |X|$; voir
Properties of Spaces,
Lemma \ref{spaces-properties-lemma-etale-image-open}.
Les morphismes $U_i \to X_i$ sont surjectifs.
Ainsi, $U_i \to X_i$ est étale surjectif, et les projections
$U_i \times_{X_i} U_i \to U_i$ sont quasi-compactes, car
$U_i \times_{X_i} U_i = U_i \times_X U_i$. Par conséquent, d'après
Spaces, Lemma \ref{spaces-lemma-representable-morphisms-spaces-property},
les morphismes $U_i \to X_i$ sont quasi-compacts.
\end{proof}








\section{Conditions sur les espaces algébriques}
\label{section-conditions}

\noindent
Dans cette section, nous étudions les relations entre diverses conditions
naturelles sur les espaces algébriques rencontrées plus haut. Veuillez lire
la Section \ref{section-reasonable-decent}
pour vous faire une idée du sens de ces conditions.

\begin{lemma}
\label{lemma-bounded-fibres}
Soit $S$ un schéma. Soit $X$ un espace algébrique sur $S$.
Considérons les conditions suivantes sur $X$ :
\begin{itemize}
\item[] $(\alpha)$ Pour tout $x \in |X|$, les conditions équivalentes du
Lemme \ref{lemma-U-finite-above-x}
sont vérifiées.
\item[] $(\beta)$ Pour tout $x \in |X|$, les conditions équivalentes du
Lemme \ref{lemma-R-finite-above-x}
sont vérifiées.
\item[] $(\gamma)$ Pour tout $x \in |X|$, les conditions équivalentes du
Lemme \ref{lemma-UR-finite-above-x}
sont vérifiées.
\item[] $(\delta)$ Les conditions équivalentes du
Lemme \ref{lemma-U-universally-bounded}
sont vérifiées.
\item[] $(\epsilon)$ Les conditions équivalentes du
Lemme \ref{lemma-characterize-very-reasonable}
sont vérifiées.
\item[] $(\zeta)$ L'espace $X$ est localement quasi-séparé pour la topologie
de Zariski.
\item[] $(\eta)$ L'espace $X$ est quasi-séparé.
\item[] $(\theta)$ L'espace $X$ est représentable, c'est-à-dire que $X$ est un
schéma.
\item[] $(\iota)$ L'espace $X$ est un schéma quasi-séparé.
\end{itemize}
Nous avons
$$
\xymatrix{
& (\theta) \ar@{=>}[rd] & & & &  \\
(\iota) \ar@{=>}[ru] \ar@{=>}[rd] & &
(\zeta) \ar@{=>}[r] &
(\epsilon) \ar@{=>}[r] &
(\delta) \ar@{=>}[r] &
(\gamma) \ar@{<=>}[r] & (\alpha) + (\beta) \\
& (\eta) \ar@{=>}[ru] & & & &
}
$$
\end{lemma}

\begin{proof}
L'implication $(\gamma) \Leftrightarrow (\alpha) + (\beta)$ est immédiate.
Les implications du losange de gauche résultent clairement des définitions.

\medskip\noindent
Supposons $(\zeta)$, c'est-à-dire que $X$ soit localement quasi-séparé pour
la topologie de Zariski. Alors $(\epsilon)$ est vérifiée d'après
Properties of Spaces, Lemma
\ref{spaces-properties-lemma-quasi-separated-quasi-compact-pieces}.

\medskip\noindent
Supposons $(\epsilon)$. D'après le
Lemme \ref{lemma-characterize-very-reasonable},
il existe un recouvrement ouvert de Zariski $X = \bigcup X_i$ tel que, pour
chaque $i$, il existe un schéma $U_i$ et un morphisme étale surjectif
quasi-compact $U_i \to X_i$. Fixons un $i$ et un sous-schéma ouvert affine
$W \subset U_i$. Il suffit de montrer que $W \to X$ a des fibres
universellement bornées, puisque la famille de tous ces morphismes
$W \to X$ recouvre alors $X$.
Pour cela, considérons le diagramme
$$
\xymatrix{
W \times_X U_i \ar[r]_-p \ar[d]_q & U_i \ar[d] \\
W \ar[r] & X
}
$$
Puisque $W \to X$ se factorise par $X_i$, nous voyons que
$W \times_X U_i = W \times_{X_i} U_i$; par conséquent, $q$ est quasi-compact.
Comme $W$ est affine, il en résulte que le schéma $W \times_X U_i$
est quasi-compact. Nous pouvons donc appliquer
Morphisms, Lemma
\ref{morphisms-lemma-locally-quasi-finite-qc-source-universally-bounded},
et nous concluons que $p$ a des fibres universellement bornées. Du
Lemme \ref{lemma-descent-universally-bounded},
nous concluons que $W \to X$ a lui aussi des fibres universellement bornées.

\medskip\noindent
Supposons $(\delta)$. Soit $U$ un schéma affine, et soit $U \to X$ un morphisme
étale. Par hypothèse, les fibres du morphisme $U \to X$ sont universellement
bornées. Il en va donc de même des fibres des deux projections
$R = U \times_X U \to U$; voir le
Lemme \ref{lemma-base-change-universally-bounded}.
Et, d'après le
Lemme \ref{lemma-composition-universally-bounded},
les fibres de $R \to X$ sont elles aussi universellement bornées.
Par conséquent, pour tout $x \in X$, les fibres de $|U| \to |X|$ et de
$|R| \to |X|$ au-dessus de $x$ sont finies; voir le
Lemme \ref{lemma-universally-bounded-finite-fibres}.
Autrement dit, les conditions équivalentes du
Lemme \ref{lemma-UR-finite-above-x}
sont vérifiées. Cela prouve $(\delta) \Rightarrow (\gamma)$.
\end{proof}

\begin{lemma}
\label{lemma-properties-local}
Soit $S$ un schéma.
Soit $\mathcal{P}$ l'une des propriétés
$(\alpha)$, $(\beta)$, $(\gamma)$, $(\delta)$, $(\epsilon)$, $(\zeta)$ ou
$(\theta)$ des espaces algébriques énumérées dans le
Lemme \ref{lemma-bounded-fibres}.
Alors, si $X$ est un espace algébrique sur $S$ et si
$X = \bigcup X_i$ est un recouvrement ouvert de Zariski tel que chaque
$X_i$ possède $\mathcal{P}$, alors $X$ possède $\mathcal{P}$.
\end{lemma}

\begin{proof}
Soit $X$ un espace algébrique sur $S$, et soit $X = \bigcup X_i$ un
recouvrement ouvert de Zariski tel que chaque $X_i$ possède $\mathcal{P}$.

\medskip\noindent
Le cas $\mathcal{P} = (\alpha)$. La condition $(\alpha)$ pour $X_i$
signifie que, pour tout $x \in |X_i|$, tout schéma affine $U$ et tout
morphisme étale $\varphi : U \to X_i$, la fibre de
$\varphi : |U| \to |X_i|$ au-dessus de $x$ est finie. Considérons
$x \in X$, un schéma affine $U$ et un morphisme étale $U \to X$.
Puisque $X = \bigcup X_i$ est un recouvrement ouvert de Zariski, il existe
un recouvrement ouvert affine fini $U = U_1 \cup \ldots \cup U_n$ tel que
chaque $U_j \to X$ se factorise par un certain $X_{i_j}$. Par hypothèse,
les fibres de $|U_j | \to |X_{i_j}|$ au-dessus de $x$ sont finies pour
$j = 1, \ldots, n$. Cela signifie clairement que la fibre de
$|U| \to |X|$ au-dessus de $x$ est finie.
Ceci prouve le résultat pour $(\alpha)$.

\medskip\noindent
Le cas $\mathcal{P} = (\beta)$. La condition $(\beta)$ pour $X_i$ signifie
que tout $x \in |X_i|$ est représenté par un monomorphisme du spectre
d'un corps vers $X_i$. Il en va donc de même pour $X$, puisque
$X_i \to X$ est un monomorphisme et que $X = \bigcup X_i$.

\medskip\noindent
Le cas $\mathcal{P} = (\gamma)$.
Notons que $(\gamma) = (\alpha) + (\beta)$ d'après le
Lemme \ref{lemma-bounded-fibres}; le lemme pour $(\gamma)$ résulte donc
des cas traités ci-dessus.

\medskip\noindent
Le cas $\mathcal{P} = (\delta)$. La condition $(\delta)$ pour $X_i$ signifie
qu'il existe des schémas $U_{ij}$ et des morphismes étales
$U_{ij} \to X_i$ à fibres universellement bornées qui recouvrent $X_i$.
Ces schémas fournissent aussi un morphisme étale surjectif
$\coprod U_{ij} \to X$, et $U_{ij} \to X$ a encore des fibres
universellement bornées.

\medskip\noindent
Le cas $\mathcal{P} = (\epsilon)$. La condition $(\epsilon)$ pour $X_i$
signifie que l'on peut trouver un ensemble $J_i$ et des morphismes
$\varphi_{ij} : U_{ij} \to X_i$ tels que chaque $\varphi_{ij}$ soit étale,
que les deux projections $U_{ij} \times_{X_i} U_{ij} \to U_{ij}$ soient
quasi-compactes et que $\coprod_{j \in J_i} U_{ij} \to X_i$ soit surjectif.
Dans ce cas, les composés $U_{ij} \to X_i \to X$ sont étales
(combiner
Morphisms, Lemmas
\ref{morphisms-lemma-composition-etale} et
\ref{morphisms-lemma-open-immersion-etale}
et
Spaces, Lemmas
\ref{spaces-lemma-composition-representable-transformations-property} et
\ref{spaces-lemma-morphism-schemes-gives-representable-transformation-property}
).
Comme $X_i \subset X$ est un sous-espace, nous voyons que
$U_{ij} \times_{X_i} U_{ij} = U_{ij} \times_X U_{ij}$; la condition sur
les produits fibrés est donc préservée. Et il est clair que
$\coprod_{i, j} U_{ij} \to X$ est surjectif. Ainsi, $X$
satisfait $(\epsilon)$.

\medskip\noindent
Le cas $\mathcal{P} = (\zeta)$. La condition $(\zeta)$ pour $X_i$ signifie
que $X_i$ est localement quasi-séparé pour la topologie de Zariski.
Il est immédiat que $X$ est alors localement quasi-séparé pour la
topologie de Zariski.

\medskip\noindent
Pour $(\theta)$, voir
Properties of Spaces,
Lemma \ref{spaces-properties-lemma-subscheme}.
\end{proof}

\begin{lemma}
\label{lemma-representable-properties}
Soit $S$ un schéma. Soit $\mathcal{P}$ l'une des propriétés
$(\beta)$, $(\gamma)$, $(\delta)$, $(\epsilon)$ ou
$(\theta)$ des espaces algébriques énumérées dans le
Lemme \ref{lemma-bounded-fibres}.
Soient $X$, $Y$ des espaces algébriques sur $S$.
Soit $X \to Y$ un morphisme représentable.
Si $Y$ possède la propriété $\mathcal{P}$, il en est de même de $X$.
\end{lemma}

\begin{proof}
Supposons que $f : X \to Y$ soit un morphisme représentable d'espaces
algébriques et que $Y$ possède $\mathcal{P}$. Soit $x \in |X|$, et posons
$y = f(x) \in |Y|$.

\medskip\noindent
Le cas $\mathcal{P} = (\beta)$. La condition $(\beta)$ pour $Y$ signifie
qu'il existe un monomorphisme $\Spec(k) \to Y$ représentant $y$.
Le produit fibré $X_y = \Spec(k) \times_Y X$ est un schéma, et
$x$ correspond à un point de $X_y$, c'est-à-dire à un monomorphisme
$\Spec(k') \to X_y$. Comme $X_y \to X$ est lui aussi un monomorphisme,
nous voyons que $x$ est représenté par le monomorphisme
$\Spec(k') \to X_y \to X$.
Autrement dit, $(\beta)$ est vérifiée pour $X$.

\medskip\noindent
Le cas $\mathcal{P} = (\gamma)$. Puisque $(\gamma) \Rightarrow (\beta)$,
nous avons vu au paragraphe précédent que $y$ et $x$ peuvent être
représentés par des monomorphismes comme dans le diagramme suivant
$$
\xymatrix{
\Spec(k') \ar[r]_-x \ar[d] & X \ar[d] \\
\Spec(k) \ar[r]^-y & Y
}
$$
De plus, par définition de la propriété $(\gamma)$ au moyen du
Lemme \ref{lemma-UR-finite-above-x} (2),
il existe des schémas $V_i$ et des morphismes étales $V_i \to Y$ tels que
$\coprod V_i \to Y$ soit surjectif et que, pour chaque $i$, en posant
$R_i = V_i \times_Y V_i$, les fibres des deux applications
$$
|V_i| \longrightarrow |Y|
\quad\text{and}\quad
|R_i| \longrightarrow |Y|
$$
au-dessus de $y$ soient finies. Cela signifie que les schémas
$(V_i)_y$ et $(R_i)_y$ sont des schémas finis sur $y = \Spec(k)$.
Comme $X \to Y$ est représentable, les produits fibrés
$U_i = V_i \times_Y X$ sont des schémas. Les morphismes $U_i \to X$
sont étales, et $\coprod U_i \to X$ est surjectif. Enfin, pour chaque $i$,
nous avons
$$
(U_i)_x =
(V_i \times_Y X)_x =
(V_i)_y \times_{\Spec(k)} \Spec(k')
$$
et
$$
(U_i \times_X U_i)_x =
\left((V_i \times_Y X) \times_X (V_i \times_Y X)\right)_x =
(R_i)_y \times_{\Spec(k)} \Spec(k')
$$
de sorte que ces schémas sont finis sur $k'$, comme changements de base des
schémas finis $(V_i)_y$ et $(R_i)_y$. Il s'ensuit que $(\gamma)$ est
vérifiée pour $X$, de nouveau par la seconde condition du
Lemme \ref{lemma-UR-finite-above-x}.

\medskip\noindent
Le cas $\mathcal{P} = (\delta)$. Soit $V \to Y$ un morphisme étale, où
$V$ est un schéma affine. Puisque $Y$ possède la propriété $(\delta)$,
ce morphisme a des fibres universellement bornées. D'après le
Lemme \ref{lemma-base-change-universally-bounded},
le changement de base $V \times_Y X \to X$ a lui aussi des fibres
universellement bornées. La première partie du
Lemme \ref{lemma-U-universally-bounded}
s'applique donc, et nous voyons que $X$ possède lui aussi la propriété
$(\delta)$.

\medskip\noindent
Le cas $\mathcal{P} = (\epsilon)$. Nous utiliserons à plusieurs reprises
Spaces, Lemma
\ref{spaces-lemma-base-change-representable-transformations-property}.
Soient $V_i \to Y$ comme dans le
Lemme \ref{lemma-characterize-very-reasonable} (2).
Posons $U_i = X \times_Y V_i$. Les morphismes $U_i \to X$ sont étales,
et $\coprod U_i \to X$ est surjectif. Puisque
$U_i \times_X U_i = X \times_Y (V_i \times_Y V_i)$, nous voyons que les
projections $U_i \times_X U_i \to U_i$ sont les changements de base des
projections $V_i \times_Y V_i \to V_i$, et sont donc elles aussi
quasi-compactes. Ainsi, $X$ satisfait la condition (2) du
Lemme \ref{lemma-characterize-very-reasonable}.

\medskip\noindent
Le cas $\mathcal{P} = (\theta)$. Dans ce cas, le résultat est
Categories, Lemma \ref{categories-lemma-representable-over-representable}.
\end{proof}











\section{Espaces algébriques raisonnables et décents}
\label{section-reasonable-decent}

\noindent
Dans le
Lemme \ref{lemma-bounded-fibres},
nous avons rencontré plusieurs conditions sur les espaces algébriques liées
au comportement des morphismes étales de schémas affines vers $X$
et à l'existence de recouvrements étales particuliers de $X$ par des
schémas. Nous récapitulons ici les différents types de conditions :
$$
\boxed{
\begin{matrix}
(\alpha) & \text{les fibres des morphismes étales issus d'affines
sont finies} \\
(\beta) & \text{les points proviennent de monomorphismes
de spectres de corps} \\
(\gamma) & \text{les points proviennent de monomorphismes quasi-compacts de
spectres de corps} \\
(\delta) & \text{les fibres des morphismes étales issus d'affines sont
universellement bornées} \\
(\epsilon) & \text{recouvrement par des morphismes étales issus de schémas,
quasi-compacts sur leur image}
\end{matrix}
}
$$

\medskip\noindent
Les conditions de la définition suivante ne sont pas exactement des conditions
sur la diagonale de $X$, mais ce sont, en un certain sens, des conditions de
séparation sur $X$.

\begin{definition}
\label{definition-very-reasonable}
Soit $S$ un schéma.
Soit $X$ un espace algébrique sur $S$.
\begin{enumerate}
\item Nous disons que $X$ est {\it décent} si, pour tout point $x \in X$, les
conditions équivalentes du
Lemme \ref{lemma-UR-finite-above-x}
sont vérifiées, autrement dit si la propriété $(\gamma)$ du
Lemme \ref{lemma-bounded-fibres}
est vérifiée.
\item Nous disons que $X$ est {\it raisonnable} si les conditions équivalentes
du
Lemme \ref{lemma-U-universally-bounded}
sont vérifiées, autrement dit si la propriété $(\delta)$ du
Lemme \ref{lemma-bounded-fibres}
est vérifiée.
\item Nous disons que $X$ est {\it très raisonnable} si les conditions
équivalentes du
Lemme \ref{lemma-characterize-very-reasonable}
sont vérifiées, c'est-à-dire si la propriété $(\epsilon)$ du
Lemme \ref{lemma-bounded-fibres}
est vérifiée.
\end{enumerate}
\end{definition}

\noindent
Nous avons les implications suivantes entre ces conditions sur les espaces
algébriques :
$$
\xymatrix{
\text{représentable} \ar@{=>}[rd] & & & \\
 & \text{très raisonnable} \ar@{=>}[r] &
\text{raisonnable} \ar@{=>}[r] &
\text{décent} \\
\text{quasi-séparé} \ar@{=>}[ru] & & &
}
$$
La notion d'espace algébrique très raisonnable est obsolète.
Elle a été introduite parce que cette hypothèse était nécessaire pour
démontrer certains résultats qui le sont maintenant pour la classe des
espaces décents.
La classe des espaces décents est la plus grande classe d'espaces $X$ pour
lesquels on dispose d'une bonne relation entre la topologie de $|X|$ et les
propriétés de $X$ lui-même.

\begin{example}
\label{example-not-decent}
L'espace algébrique $\mathbf{A}^1_{\mathbf{Q}}/\mathbf{Z}$ construit dans
Spaces, Example \ref{spaces-example-affine-line-translation}
n'est pas décent, car son ``point générique'' ne peut pas être représenté par
un monomorphisme du spectre d'un corps.
\end{example}

\begin{remark}
\label{remark-reasonable}
Les espaces algébriques raisonnables sont techniquement plus faciles à manier
que les espaces algébriques très raisonnables. Par exemple, si $X \to Y$ est
un morphisme quasi-compact, étale et surjectif d'espaces algébriques et si
$X$ est raisonnable, alors $Y$ l'est aussi; voir le
Lemme \ref{lemma-descent-conditions}.
Mais nous ne savons pas si cela reste vrai pour la propriété
``très raisonnable''. Nous donnons plus bas une autre propriété technique
que possèdent les espaces algébriques raisonnables.
\end{remark}

\begin{lemma}
\label{lemma-fun-property-reasonable}
Soit $S$ un schéma.
Soit $X$ un espace algébrique raisonnable quasi-compact sur $S$.
Il existe alors un système inductif filtrant d'espaces algébriques
quasi-compacts et quasi-séparés $X_i$ tel que $X = \colim_i X_i$
(la colimite étant prise dans la catégorie des faisceaux). De plus, on peut
faire en sorte que
\begin{enumerate}
\item pour tout schéma quasi-compact $T$ sur $S$, on ait
$\colim X_i(T) = X(T)$,
\item les morphismes de transition $X_i \to X_{i'}$ du système et les
coprojections $X_i \to X$ soient surjectifs et étales, et
\item si $X$ est un schéma, les espaces algébriques $X_i$ soient des schémas,
et les morphismes de transition $X_i \to X_{i'}$ et les coprojections
$X_i \to X$ soient des isomorphismes locaux.
\end{enumerate}
\end{lemma}

\begin{proof}
Esquissons la démonstration. D'après
Properties of Spaces, Lemma
\ref{spaces-properties-lemma-quasi-compact-affine-cover},
nous avons $X = U/R$, avec $U$ affine.
Dans ce cas, être raisonnable signifie que $U \to X$ est
universellement borné. Il existe donc un entier $N$ tel que les ``fibres'' de
$U \to X$ soient de degré au plus $N$; voir la
Définition \ref{definition-universally-bounded}.
Désignons par $s, t : R \to U$ et $c : R \times_{s, U, t} R \to R$ les
applications structurales du groupoïde.

\medskip\noindent
Affirmation : pour tout ouvert quasi-compact $A \subset R$, il existe un ouvert
$R' \subset R$ tel que
\begin{enumerate}
\item $A \subset R'$,
\item $R'$ soit quasi-compact, et
\item $(U, R', s|_{R'}, t|_{R'}, c|_{R' \times_{s, U, t} R'})$ soit
un groupoïde en schémas.
\end{enumerate}
Notons que le morphisme $e : U \to R$ est ouvert, puisqu'il est une section du
morphisme étale $s : R \to U$; voir
\'Etale Morphisms, Proposition \ref{etale-proposition-properties-sections}.
De plus, $U$ est affine, donc quasi-compact. Nous pouvons ainsi remplacer $A$
par $A \cup e(U) \subset R$ et supposer que $A$ contient $e(U)$. Définissons
ensuite par récurrence $A^1 = A$ et
$$
A^n = c(A^{n - 1} \times_{s, U, t} A) \subset R
$$
pour $n \geq 2$. Un raisonnement par récurrence montre que $A^n$ est
quasi-compact pour tout $n \geq 2$, comme image du produit fibré quasi-compact
$A^{n - 1} \times_{s, U, t} A$. Si $k$ est un corps algébriquement clos sur
$S$ et si nous considérons les points à valeurs dans $k$, alors
$$
A^n(k) = \left\{(u, u') \in U(k) \times U(k)
:
\begin{matrix}
\text{il existe } u = u_1, u_2, \ldots, u_{n + 1} = u' \in U(k)
\text{ tels que} \\
(u_i , u_{i + 1}) \in A \text{ pour tout }i = 1, \ldots, n.
\end{matrix}
\right\}
$$
Mais, comme les fibres de $U(k) \to X(k)$ sont de cardinal au plus $N$, dès
que $n > N$, un terme se répète dans la suite ci-dessus et nous pouvons la
raccourcir. Cela prouve que $A^N = A^n$ pour tout $n \geq N$.
Il en résulte que $R' = A^N$ fournit le groupoïde en schémas
$(U, R', s|_{R'}, t|_{R'}, c|_{R' \times_{s, U, t} R'})$, ce qui démontre
l'affirmation.

\medskip\noindent
Considérons le morphisme de faisceaux sur $(\Sch/S)_{fppf}$
$$
\colim_{R' \subset R} U/R' \longrightarrow U/R
$$
où $R' \subset R$ parcourt les sous-schémas ouverts quasi-compacts de $R$
qui définissent des relations d'équivalence étales comme ci-dessus.
Chacun des quotients $U/R'$ est un espace algébrique
(voir Spaces, Theorem \ref{spaces-theorem-presentation}).
Comme $R'$ est quasi-compact et $U$ affine, le morphisme
$R' \to U \times_{\Spec(\mathbf{Z})} U$ est quasi-compact;
par conséquent, $U/R'$ est quasi-séparé. Enfin, si $T$ est un schéma
quasi-compact, alors
$$
\colim_{R' \subset R} U(T)/R'(T) \longrightarrow U(T)/R(T)
$$
est une bijection, car tout morphisme de $T$ dans $R$ se factorise par l'une
des sous-relations ouvertes $R'$ grâce à l'affirmation ci-dessus. Cela
implique clairement que la colimite des faisceaux $U/R'$ est $U/R$. Autrement
dit, l'espace algébrique $X = U/R$ est la colimite des espaces algébriques
quasi-séparés $U/R'$.

\medskip\noindent
Les propriétés (1) et (2) résultent de la discussion précédente.
Si $X$ est un schéma, choisir pour $U$ une réunion disjointe finie d'ouverts
affines de $X$ donne (3).
Nous omettons les détails.
\end{proof}

\begin{lemma}
\label{lemma-representable-named-properties}
Soit $S$ un schéma. Soient $X$, $Y$ des espaces algébriques sur $S$.
Soit $X \to Y$ un morphisme représentable.
Si $Y$ est décent (resp.\ raisonnable), alors $X$ l'est aussi.
\end{lemma}

\begin{proof}
C'est une reformulation du Lemme \ref{lemma-representable-properties}.
\end{proof}

\begin{lemma}
\label{lemma-etale-named-properties}
Soit $S$ un schéma. Soit $X \to Y$ un morphisme étale d'espaces
algébriques sur $S$. Si $Y$ est décent, resp.\ raisonnable,
alors $X$ l'est aussi.
\end{lemma}

\begin{proof}
Soit $U$ un schéma affine et $U \to X$ un morphisme étale.
Posons $R = U \times_X U$ et $R' = U \times_Y U$. Notons que
$R \to R'$ est un monomorphisme.

\medskip\noindent
Soit $x \in |X|$. Pour montrer que $X$ est décent, nous devons montrer que
les fibres de $|U| \to |X|$ et de $|R| \to |X|$ au-dessus de $x$ sont finies.
Mais si $Y$ est décent, les fibres de $|U| \to |Y|$ et de
$|R'| \to |Y|$ sont finies. D'où le résultat pour ``décent''.

\medskip\noindent
Pour montrer que $X$ est raisonnable, nous devons montrer que les fibres de
$U \to X$ sont universellement bornées. Or, si $Y$ est raisonnable, les fibres
de $U \to Y$ sont universellement bornées, ce qui implique immédiatement la
même propriété pour les fibres de $U \to X$.
D'où le résultat pour ``raisonnable''.
\end{proof}







\section{Points et spécialisations}
\label{section-specializations}

\noindent
Il existe un morphisme étale d'espaces algébriques $f : X \to Y$ et une
spécialisation non triviale entre des points d'une fibre de
$|f| : |X| \to |Y|$; voir
Examples, Lemma \ref{examples-lemma-specializations-fibre-etale}.
Si la source du morphisme est un schéma, on peut exclure ce phénomène en
imposant la condition ($\alpha$) à $Y$.

\begin{lemma}
\label{lemma-no-specializations-map-to-same-point}
Soit $S$ un schéma.
Soit $X$ un espace algébrique sur $S$.
Soit $U \to X$ un morphisme étale d'un schéma vers $X$.
Supposons que $u, u' \in |U|$ aient pour image le même point $x$ de $|X|$ et
que $u' \leadsto u$. Si le couple $(X, x)$ satisfait les conditions
équivalentes du
Lemme \ref{lemma-U-finite-above-x},
alors $u = u'$.
\end{lemma}

\begin{proof}
Supposons que le couple $(X, x)$ satisfasse les conditions équivalentes du
Lemme \ref{lemma-U-finite-above-x}.
Soit $U$ un schéma, soit $U \to X$ un morphisme étale, et supposons que
$u, u' \in |U|$ aient pour image $x$ dans $|X|$ et que
$u' \leadsto u$. Nous pouvons remplacer, et remplaçons, $U$ par un voisinage
affine de $u$. Soient $t, s : R = U \times_X U \to U$
les projections étales.

\medskip\noindent
Choisissons un point $r \in R$ tel que $t(r) = u$ et $s(r) = u'$.
C'est possible d'après
Properties of Spaces,
Lemma \ref{spaces-properties-lemma-points-presentation}.
Comme les générisations se relèvent le long du morphisme étale $t$
(Remarque \ref{remark-recall}), nous pouvons trouver une spécialisation
$r' \leadsto r$ telle que $t(r') = u'$. Posons $u'' = s(r')$. Alors
$u'' \leadsto u'$.
Nous pouvons donc recommencer et trouver $r'' \leadsto r'$ tel que
$t(r'') = u''$. Posons $u''' = s(r'')$, et ainsi de suite.
Voici le diagramme :
$$
\xymatrix{
& r'' \ar[rd]^s \ar[ld]_t \ar@{~>}[d] & \\
u'' \ar@{~>}[d] & r' \ar[rd]^s \ar[ld]_t \ar@{~>}[d] & u''' \ar@{~>}[d] \\
u' \ar@{~>}[d] & r \ar[rd]^s \ar[ld]_t & u'' \ar@{~>}[d] \\
u & & u'
}
$$
Dans la Remarque \ref{remark-recall}, nous avons vu qu'il n'y a pas de
spécialisations entre les points des fibres du morphisme étale $s$. Ainsi,
si $u^{(n + 1)} = u^{(n)}$ pour un certain $n$, alors aussi
$r^{(n)} = r^{(n - 1)}$, puis, en appliquant $t$,
$u^{(n)} = u^{(n - 1)}$. Toute la tour est alors constante; en particulier,
$u = u'$. Nous voyons donc que, si $u \not = u'$, toutes les spécialisations
sont strictes et que $\{u, u', u'', \ldots\}$ est un ensemble infini de points
de $U$ dont l'image est le point $x$ de $|X|$. Comme nous avons choisi $U$
affine, cela contredit la seconde partie du
Lemme \ref{lemma-U-finite-above-x}, comme voulu.
\end{proof}

\begin{lemma}
\label{lemma-no-specializations-map-to-same-point-Noetherian}
Soit $S$ un schéma. Soit $X$ un espace algébrique sur $S$.
Soit $U \to X$ un morphisme étale d'un schéma vers $X$.
Supposons que $u, u' \in |U|$ aient pour image le même point $x$ de $|X|$ et
que $u' \leadsto u$. Si $X$ est localement noethérien, alors $u = u'$.
\end{lemma}

\begin{proof}
La discussion de Schemes, Section \ref{schemes-section-points}
montre que $\mathcal{O}_{U, u'}$ est un localisé de
l'anneau local noethérien $\mathcal{O}_{U, u}$.
D'après Properties of Spaces, Lemma
\ref{spaces-properties-lemma-pre-dimension-local-ring},
nous avons $\dim(\mathcal{O}_{U, u}) = \dim(\mathcal{O}_{U, u'})$.
La théorie de la dimension des anneaux locaux noethériens permet de conclure
que $u = u'$.
\end{proof}

\begin{lemma}
\label{lemma-specialization}
Soit $S$ un schéma.
Soit $X$ un espace algébrique sur $S$.
Soient $x, x' \in |X|$ et supposons que $x' \leadsto x$, c'est-à-dire que $x$
est une spécialisation de $x'$.
Supposons que le couple $(X, x')$ satisfasse les conditions équivalentes
du Lemme \ref{lemma-UR-finite-above-x}.
Alors, pour tout morphisme étale $\varphi : U \to X$ d'un schéma $U$ et tout
$u \in U$ tel que $\varphi(u) = x$, il existe un point $u'\in U$ tel que
$u' \leadsto u$ et $\varphi(u') = x'$.
\end{lemma}

\begin{proof}
Nous pouvons remplacer $U$ par un voisinage ouvert affine de $u$.
Nous pouvons donc supposer que $U$ est affine. Comme $x$ appartient à l'image
du morphisme ouvert $|U| \to |X|$, il en est de même de $x'$. Nous pouvons
ainsi remplacer $X$ par le sous-espace ouvert de Zariski correspondant à
l'image de $|U| \to |X|$; voir
Properties of Spaces,
Lemma \ref{spaces-properties-lemma-etale-image-open}.
Autrement dit, nous pouvons supposer que
$U \to X$ est surjectif et étale.
Soient $s, t : R = U \times_X U \to U$ les projections.
Comme $(X, x')$ satisfait par hypothèse les conditions équivalentes du
Lemme \ref{lemma-UR-finite-above-x},
les fibres de $|U| \to |X|$ et de $|R| \to |X|$
au-dessus de $x'$ sont finies. Écrivons $\{u'_1, \ldots, u'_n\} \subset U$ et
$\{r'_1, \ldots, r'_m\} \subset R$ pour les images inverses respectives
de $\{x'\}$.
Considérons les fermés
$$
T = \overline{\{u'_1\}} \cup \ldots \cup \overline{\{u'_n\}} \subset |U|,
\quad
T' = \overline{\{r'_1\}} \cup \ldots \cup \overline{\{r'_m\}} \subset |R|.
$$
Il est clair que $s(T') \subset T$. Comme $R$ est une relation d'équivalence,
nous avons aussi $t(T') = s(T')$, puisque l'ensemble $\{r_j'\}$ est invariant
par la symétrie de $R$, par construction. Soit $w \in T$ un point quelconque.
Alors $u'_i \leadsto w$ pour un certain $i$. Choisissons $r \in R$
tel que $s(r) = w$. Comme les générisations se relèvent le long de
$s : R \to U$ (voir la Remarque \ref{remark-recall}), nous pouvons trouver
$r' \leadsto r$ tel que $s(r') = u_i'$. Alors $r' = r'_j$ pour un certain $j$,
et nous en déduisons que $w \in s(T')$. Ainsi,
$T = s(T') = t(T')$ est une partie saturée pour la relation d'équivalence
$|R|$, fermée dans $|U|$. Cela signifie que $T$ est l'image inverse d'une
partie fermée (!) $T'' = \varphi(T)$ de $|X|$; voir
Properties of Spaces,
Lemmas \ref{spaces-properties-lemma-points-presentation} et
\ref{spaces-properties-lemma-topology-points}.
Par conséquent, $T'' = \overline{\{x'\}}$.
Ainsi, $T$ contient un point $u_1$ d'image $x$, puisque $x \in T''$.
Autrement dit, pour un certain $i$, il existe une spécialisation
$u'_i \leadsto u_1$ dont l'image est la spécialisation donnée
$x' \leadsto x$.

\medskip\noindent
Pour achever la preuve, choisissons un point $r \in R$ tel que
$s(r) = u$ et $t(r) = u_1$ (à l'aide de
Properties of Spaces,
Lemma \ref{spaces-properties-lemma-points-cartesian}).
Comme les générisations se relèvent le long de $t$ et que
$u'_i \leadsto u_1$, nous pouvons trouver une spécialisation
$r' \leadsto r$ telle que $t(r') = u'_i$.
Posons $u' = s(r')$. Alors $u' \leadsto u$ et $\varphi(u') = x'$, comme voulu.
\end{proof}

\begin{lemma}
\label{lemma-generalizations-lift-flat}
Soit $S$ un schéma. Soit $f : Y \to X$ un morphisme plat d'espaces algébriques
sur $S$. Soient $x, x' \in |X|$ et supposons que $x' \leadsto x$, c'est-à-dire
que $x$ est une spécialisation de $x'$. Supposons que le couple $(X, x')$
satisfasse les conditions équivalentes du
Lemme \ref{lemma-UR-finite-above-x} (par exemple si
$X$ est décent, $X$ est quasi-séparé ou $X$ est représentable).
Alors, pour tout $y \in |Y|$ tel que $f(y) = x$, il existe un point
$y' \in |Y|$ tel que $y' \leadsto y$ et $f(y') = x'$.
\end{lemma}

\begin{proof}
(L'assertion entre parenthèses résulte de la définition des espaces décents
et des implications entre les différentes conditions de séparation
mentionnées dans la Section \ref{section-reasonable-decent}.)
Choisissons un schéma $V$ et un morphisme étale surjectif $V \to Y$.
Choisissons $v \in V$ d'image $y$. Il suffit alors de démontrer le lemme pour
$V \to X$. Nous pouvons donc supposer que $Y$ est un schéma.
Choisissons un schéma $U$ et un morphisme étale surjectif $U \to X$.
Choisissons $u \in U$ d'image $x$. D'après le
Lemme \ref{lemma-specialization}, nous pouvons choisir $u' \leadsto u$
d'image $x'$. D'après Properties of Spaces,
Lemma \ref{spaces-properties-lemma-points-cartesian},
nous pouvons choisir $z \in U \times_X Y$ d'image $y$ et $u$.
Nous nous ramenons ainsi au cas du morphisme plat de schémas
$U \times_X Y \to U$. La conclusion résulte alors de
Morphisms, Lemma \ref{morphisms-lemma-generalizations-lift-flat}.
\end{proof}






\section{Stratification des espaces algébriques par des schémas}
\label{section-stratifications}

\noindent
Dans cette section, nous montrons qu'un espace algébrique quasi-compact et
quasi-séparé admet une stratification finie par des sous-espaces localement
fermés dont chacun est un schéma, et telle que le recollement des strates
s'effectue par des carrés distingués élémentaires. Nous établissons d'abord
un résultat un peu plus faible pour les espaces algébriques raisonnables.

\begin{lemma}
\label{lemma-quasi-compact-reasonable-stratification}
Soit $S$ un schéma. Soit $W \to X$ un morphisme d'un schéma $W$
vers un espace algébrique $X$, plat, localement de présentation finie,
séparé, localement quasi-fini et à fibres universellement bornées. Il existe
des sous-espaces fermés réduits
$$
\emptyset = Z_{-1} \subset Z_0 \subset Z_1 \subset Z_2 \subset
\ldots \subset Z_n = X
$$
tels que, si $X_r = Z_r \setminus Z_{r - 1}$, la stratification
$X = \coprod_{r = 0, \ldots, n} X_r$ est caractérisée par la propriété
universelle suivante : étant donné $g : T \to X$, la projection
$W \times_X T \to T$ est finie localement libre de degré $r$ si et seulement si
$g(|T|) \subset |X_r|$.
\end{lemma}

\begin{proof}
Soit $n$ un entier qui majore les degrés des fibres de $W \to X$.
Choisissons un schéma $U$ et un morphisme étale surjectif $U \to X$.
Appliquons More on Morphisms, Lemma
\ref{more-morphisms-lemma-stratify-flat-fp-lqf-universally-bounded}
à $W \times_X U \to U$. Nous obtenons des parties fermées
$$
\emptyset = Y_{-1} \subset Y_0 \subset Y_1 \subset Y_2 \subset
\ldots \subset Y_n = U
$$
caractérisées par la propriété de l'énoncé pour le morphisme
$W \times_X U \to U$. Manifestement, la formation de ces parties fermées
commute au changement de base. Posant $R = U \times_X U$, et notant
$s, t : R \to U$ les projections, nous concluons que
$$
s^{-1}(Y_r) = t^{-1}(Y_r)
$$
en tant que parties fermées de $R$. Autrement dit, les parties fermées
$Y_r \subset U$ sont saturées pour la relation $R$. Cela signifie que
$|Y_r|$ est l'image inverse d'une partie fermée $Z_r \subset |X|$.
Notons encore $Z_r \subset X$ le sous-espace algébrique muni de la structure
réduite induite; voir Properties of Spaces, Definition
\ref{spaces-properties-definition-reduced-induced-space}.

\medskip\noindent
Soit $g : T \to X$ un morphisme d'espaces algébriques. Choisissons un schéma
$V$ et un morphisme étale surjectif $V \to T$. Pour démontrer la dernière
assertion du lemme, il suffit de la démontrer pour la composée
$V \to X$ (d'après notre définition des morphismes finis localement libres;
voir Morphisms of Spaces, Section
\ref{spaces-morphisms-section-finite-locally-free}).
De même, le morphisme de schémas $W \times_X V \to V$ est fini
localement libre de degré $r$ si et seulement si le morphisme de schémas
$$
W \times_X (U \times_X V)
\longrightarrow
U \times_X V
$$
est fini localement libre de degré $r$ (voir
Descent, Lemma \ref{descent-lemma-descending-property-finite-locally-free}).
Par construction, cela se produit si et seulement si
$|U \times_X V| \to |U|$ a son image contenue dans $|Y_r|$, ce qui équivaut
à ce que $|V| \to |X|$ ait son image contenue dans $|Z_r|$.
\end{proof}

\begin{lemma}
\label{lemma-stratify-flat-fp-lqf}
Soit $S$ un schéma. Soit $W \to X$ un morphisme d'un schéma $W$
vers un espace algébrique $X$, plat, localement de présentation finie,
séparé et localement quasi-fini. Il existe alors des sous-espaces ouverts
$$
X = X_0 \supset X_1 \supset X_2 \supset \ldots
$$
tels qu'un morphisme $\Spec(k) \to X$, où $k$ est un corps,
se factorise par $X_d$ si et seulement si
$W \times_X \Spec(k)$ est de degré $\geq d$ sur $k$.
\end{lemma}

\begin{proof}
Choisissons un schéma $U$ et un morphisme étale surjectif $U \to X$.
Appliquons More on Morphisms, Lemma
\ref{more-morphisms-lemma-stratify-flat-fp-lqf}
à $W \times_X U \to U$. Nous obtenons des sous-schémas ouverts
$$
U = U_0 \supset U_1 \supset U_2 \supset \ldots
$$
caractérisés par la propriété de l'énoncé pour le morphisme
$W \times_X U \to U$. Manifestement, la formation de ces sous-schémas ouverts
commute au changement de base. Posons $R = U \times_X U$ et notons
$s, t : R \to U$ les projections. Nous concluons que
$$
s^{-1}(U_d) = t^{-1}(U_d)
$$
en tant que sous-schémas ouverts de $R$. Autrement dit, les sous-schémas
ouverts $U_d \subset U$ sont saturés pour la relation $R$. Cela signifie que
$U_d$ est l'image inverse d'un sous-espace ouvert $X_d \subset X$
(Properties of Spaces, Lemma
\ref{spaces-properties-lemma-subspaces-presentation}).
\end{proof}

\begin{lemma}
\label{lemma-filter-quasi-compact}
Soit $S$ un schéma. Soit $X$ un espace algébrique quasi-compact sur $S$.
Il existe des sous-espaces ouverts
$$
\ldots \subset U_4 \subset U_3 \subset U_2 \subset U_1 = X
$$
ayant les propriétés suivantes :
\begin{enumerate}
\item si l'on pose $T_p = U_p \setminus U_{p + 1}$ (muni de la structure
réduite induite), il existe un schéma séparé $V_p$ et un morphisme étale
surjectif $f_p : V_p \to U_p$ tels que $f_p^{-1}(T_p) \to T_p$ soit un
isomorphisme,
\item si $x \in |X|$ peut être représenté par un morphisme quasi-compact
$\Spec(k) \to X$ dont la source est le spectre d'un corps, alors
$x \in T_p$ pour un certain $p$.
\end{enumerate}
\end{lemma}

\begin{proof}
D'après Properties of Spaces, Lemma
\ref{spaces-properties-lemma-quasi-compact-affine-cover},
nous pouvons choisir un schéma affine $U$ et un morphisme étale surjectif
$U \to X$. Pour $p \geq 0$, posons
$$
W_p = U \times_X \ldots \times_X U \setminus
\text{toutes les diagonales}
$$
où le produit fibré comporte $p$ facteurs. Puisque $U$ est séparé,
le morphisme $U \to X$ est séparé et tous les produits fibrés
$U \times_X \ldots \times_X U$ sont des schémas séparés. Comme $U \to X$
est séparé, la diagonale $U \to U \times_X U$ est une immersion fermée.
Comme $U \to X$ est étale, la diagonale $U \to U \times_X U$ est une
immersion ouverte; voir Morphisms of Spaces, Lemmas
\ref{spaces-morphisms-lemma-etale-unramified} et
\ref{spaces-morphisms-lemma-diagonal-unramified-morphism}.
De même, tous les morphismes diagonaux sont des immersions ouvertes et
fermées, et $W_p$ est un sous-schéma ouvert et fermé de
$U \times_X \ldots \times_X U$. De plus, le morphisme
$$
U \times_X \ldots \times_X U \longrightarrow
U \times_{\Spec(\mathbf{Z})} \ldots \times_{\Spec(\mathbf{Z})} U
$$
est localement quasi-fini et séparé (Morphisms of Spaces,
Lemma \ref{spaces-morphisms-lemma-fibre-product-after-map}),
et son but est un schéma affine. Toute partie finie de
$U \times_X \ldots \times_X U$ est donc contenue dans un ouvert affine;
voir More on Morphisms, Lemma
\ref{more-morphisms-lemma-separated-locally-quasi-finite-over-affine}.
Il en va donc de même pour $W_p$.
Le groupe symétrique $S_p$ opère librement sur $W_p$ au-dessus de $X$
(puisque nous avons retranché le lieu des points fixes de
$U \times_X \ldots \times_X U$). D'après ce qui précède et
Properties of Spaces, Proposition
\ref{spaces-properties-proposition-finite-flat-equivalence-global},
le quotient $V_p = W_p/S_p$ est un schéma. Comme l'action de $S_p$ sur
$W_p$ est au-dessus de $X$, il existe un morphisme $V_p \to X$.
Puisque $W_p \to X$ est étale et que $W_p \to V_p$ est étale surjectif,
il s'ensuit que $V_p \to X$ est lui aussi étale; voir
Properties of Spaces, Lemma \ref{spaces-properties-lemma-etale-local}.
Notons que $V_p$ est un schéma séparé d'après
Properties of Spaces, Lemma
\ref{spaces-properties-lemma-quotient-separated}.

\medskip\noindent
Soit $U_p \subset X$ le sous-espace ouvert image de $V_p \to X$.
Par construction, un morphisme $\Spec(k) \to X$, où $k$ est
algébriquement clos, se factorise par $U_p$ si et seulement si
$U \times_X \Spec(k)$ possède $\geq p$ points; comme d'habitude, notons que
$U \times_X \Spec(k)$ est, en tant que schéma, une réunion disjointe
(éventuellement infinie) de copies de $\Spec(k)$; voir
Remark \ref{remark-recall}. Il s'ensuit que les $U_p$ forment une filtration
de $X$ comme dans l'énoncé du lemme. De plus, notre morphisme
$\Spec(k) \to X$ se factorise par $T_p$ si et seulement si
$U \times_X \Spec(k)$ possède exactement $p$ points. Dans ce cas, nous
voyons que $V_p \times_X \Spec(k)$ possède exactement un point.
Posons $Z_p = f_p^{-1}(T_p) \subset V_p$. C'est un sous-schéma fermé de
$V_p$. Alors $Z_p \to T_p$ est un morphisme étale d'espaces algébriques qui
induit une bijection sur les points à valeurs dans $k$ pour tout corps
algébriquement clos $k$. Précisément, cela implique que $Z_p \to T_p$ est
universellement injectif, donc une immersion ouverte d'après
Morphisms of Spaces, Lemma
\ref{spaces-morphisms-lemma-etale-universally-injective-open},
et par conséquent un isomorphisme. Cela démontre (1).

\medskip\noindent
Soit $x : \Spec(k) \to X$ un morphisme quasi-compact, où $k$ est un corps.
Alors la composée $\Spec(\overline{k}) \to \Spec(k) \to X$ est elle aussi
quasi-compacte (Morphisms of Spaces, Lemma
\ref{spaces-morphisms-lemma-composition-quasi-compact}).
Dans ce cas, le schéma $U \times_X \Spec(\overline{k})$ est quasi-compact.
Comme nous avons vu ci-dessus qu'il est une réunion disjointe de copies de
$\Spec(\overline{k})$, nous constatons qu'il a un nombre fini de points.
Si ce nombre est $p$, alors $x \in T_p$, comme voulu, ce qui achève la preuve.
\end{proof}

\begin{lemma}
\label{lemma-filter-reasonable}
Soit $S$ un schéma. Soit $X$ un espace algébrique quasi-compact et raisonnable
sur $S$. Il existe un entier $n$ et des sous-espaces ouverts
$$
\emptyset = U_{n + 1} \subset
U_n \subset U_{n - 1} \subset \ldots \subset U_1 = X
$$
ayant la propriété suivante : si l'on pose $T_p = U_p \setminus U_{p + 1}$
(muni de la structure réduite induite), il existe un schéma séparé $V_p$ et
un morphisme étale surjectif $f_p : V_p \to U_p$ tels que
$f_p^{-1}(T_p) \to T_p$ soit un isomorphisme.
\end{lemma}

\begin{proof}
La démonstration de ce lemme est identique à celle du
Lemme \ref{lemma-filter-quasi-compact}.
Soit $n$ un entier qui majore les degrés des fibres de $U \to X$;
un tel entier existe puisque $X$ est raisonnable, voir
Definition \ref{definition-very-reasonable}.
Nous constatons alors que $U_{n + 1} = \emptyset$, ce qui achève la preuve.
\end{proof}

\begin{lemma}
\label{lemma-stratify-reasonable}
Soit $S$ un schéma. Soit $X$ un espace algébrique quasi-compact et raisonnable
sur $S$. Il existe un entier $n$ et des sous-espaces ouverts
$$
\emptyset = U_{n + 1} \subset
U_n \subset U_{n - 1} \subset \ldots \subset U_1 = X
$$
tels que chaque $T_p = U_p \setminus U_{p + 1}$ (muni de la structure réduite
induite) soit un schéma.
\end{lemma}

\begin{proof}
Conséquence immédiate du Lemme \ref{lemma-filter-reasonable}.
\end{proof}

\noindent
Le résultat suivant est presque identique à
\cite[Proposition 5.7.8]{GruRay}.

\begin{lemma}
\label{lemma-filter-quasi-compact-quasi-separated}
\begin{reference}
Ce résultat est presque identique à \cite[Proposition 5.7.8]{GruRay}.
\end{reference}
Soit $X$ un espace algébrique quasi-compact et quasi-séparé sur
$\Spec(\mathbf{Z})$. Il existe un entier $n$ et des sous-espaces ouverts
$$
\emptyset = U_{n + 1} \subset
U_n \subset U_{n - 1} \subset \ldots \subset U_1 = X
$$
ayant la propriété suivante : si l'on pose $T_p = U_p \setminus U_{p + 1}$
(muni de la structure réduite induite), il existe un schéma $V_p$
quasi-compact et séparé et un morphisme étale surjectif
$f_p : V_p \to U_p$ tels que $f_p^{-1}(T_p) \to T_p$ soit un isomorphisme.
\end{lemma}

\begin{proof}
La démonstration de ce lemme est identique à celle du
Lemme \ref{lemma-filter-quasi-compact}.
Remarquons qu'un espace quasi-séparé est raisonnable; voir
Lemma \ref{lemma-bounded-fibres} et
Definition \ref{definition-very-reasonable}.
Nous obtenons donc $U_{n + 1} = \emptyset$, comme dans le
Lemme \ref{lemma-filter-reasonable}.
À la fin de l'argument, ajoutons que, puisque $X$ est quasi-séparé,
les schémas $U \times_X \ldots \times_X U$ sont tous quasi-compacts.
Les schémas $W_p$ sont donc quasi-compacts. Par conséquent, les quotients
$V_p = W_p/S_p$ par le groupe symétrique $S_p$ sont des schémas
quasi-compacts.
\end{proof}

\noindent
Le lemme suivant devrait probablement se trouver ailleurs.

\begin{lemma}
\label{lemma-locally-constructible}
Soit $S$ un schéma. Soit $X$ un espace algébrique quasi-séparé sur $S$.
Soit $E \subset |X|$ une partie. Alors $E$ est étale-localement constructible
(Properties of Spaces, Definition
\ref{spaces-properties-definition-locally-constructible})
si et seulement si $E$ est une partie localement constructible de
l'espace topologique $|X|$
(Topology, Definition \ref{topology-definition-constructible}).
\end{lemma}

\begin{proof}
Supposons que $E \subset |X|$ soit une partie localement constructible de
l'espace topologique $|X|$. Soit $f : U \to X$ un morphisme étale, où $U$
est un schéma. Nous devons montrer que $f^{-1}(E)$ est localement constructible
dans $U$. La question est locale sur $U$ et $X$; nous pouvons donc supposer
que $X$ est quasi-compact, que $E \subset |X|$ est constructible et que $U$
est affine. Dans ce cas, $U \to X$ est quasi-compact, donc
$f : |U| \to |X|$ est quasi-compact. Remarquons que les ouverts rétrocompacts
de $|X|$, resp.\ $U$, ne sont autres que les ouverts quasi-compacts de
$|X|$, resp.\ $U$; voir
Topology, Lemma \ref{topology-lemma-topology-quasi-separated-scheme}.
Ainsi, $f^{-1}(E)$ est constructible d'après Topology, Lemma
\ref{topology-lemma-inverse-images-constructibles}.

\medskip\noindent
Réciproquement, supposons que $E$ soit étale-localement constructible.
Nous voulons montrer que $E$ est localement constructible dans l'espace
topologique $|X|$. La question est locale sur $X$; nous pouvons donc supposer
que $X$ est à la fois quasi-compact et quasi-séparé. Montrons que, dans ce cas,
$E$ est constructible dans $|X|$. Choisissons des sous-espaces ouverts
$$
\emptyset = U_{n + 1} \subset
U_n \subset U_{n - 1} \subset \ldots \subset U_1 = X
$$
et des morphismes étales surjectifs $f_p : V_p \to U_p$ induisant des
isomorphismes $f_p^{-1}(T_p) \to T_p = U_p \setminus U_{p + 1}$, où $V_p$
est un schéma quasi-compact et séparé, comme dans le
Lemme \ref{lemma-filter-quasi-compact-quasi-separated}.
Par définition, l'image inverse $E_p \subset V_p$ de $E$ est localement
constructible dans $V_p$. Alors $E_p$ est constructible dans $V_p$ d'après
Properties, Lemma
\ref{properties-lemma-constructible-quasi-compact-quasi-separated}.
Ainsi, $E_p \cap |f_p^{-1}(T_p)| = E \cap |T_p|$ est constructible dans
$|T_p|$ d'après
Topology, Lemma \ref{topology-lemma-intersect-constructible-with-closed}
(remarquons que $V_p \setminus f_p^{-1}(T_p)$ est quasi-compact, car c'est
l'image inverse de l'espace quasi-compact $U_{p + 1}$ par le morphisme
quasi-compact $f_p$). Par conséquent,
$$
E = (|T_n| \cap E) \cup (|T_{n - 1}| \cap E) \cup \ldots \cup
(|T_1| \cap E)
$$
est constructible d'après
Topology, Lemma \ref{topology-lemma-collate-constructible-from-constructible}.
Nous utilisons ici le fait que $|T_p|$ est constructible dans $|X|$, ce qui
résulte clairement de ce qui précède.
\end{proof}





\section{Recouvrement entier par un schéma}
\label{section-integral-cover}

\noindent
Nous démontrons ici que, pour tout espace algébrique $X$ quasi-compact et
quasi-séparé, il existe un schéma $Y$ et un morphisme entier surjectif
$Y \to X$. Après avoir développé une théorie des limites d'espaces
algébriques, nous démontrerons qu'on peut obtenir un morphisme fini; voir
Limits of Spaces, Section \ref{spaces-limits-section-finite-cover}.

\begin{lemma}
\label{lemma-extend-integral-morphism}
Soit $S$ un schéma. Soit $j : V \to Y$ une immersion ouverte quasi-compacte
d'espaces algébriques sur $S$. Soit $\pi : Z \to V$ un morphisme entier.
Il existe alors un morphisme entier $\nu : Y' \to Y$ tel que $Z$ soit
$V$-isomorphe à l'image inverse de $V$ dans $Y'$.
\end{lemma}

\begin{proof}
Puisque $j$ et $\pi$ sont tous deux quasi-compacts et séparés, il en va de
même de $j \circ \pi$. Soit $\nu : Y' \to Y$ la normalisation de $Y$ dans
$Z$; voir Morphisms of Spaces, Section
\ref{spaces-morphisms-section-normalization-X-in-Y}.
Bien entendu, $\nu$ est entier; voir
Morphisms of Spaces, Lemma
\ref{spaces-morphisms-lemma-characterize-normalization}.
La dernière assertion résulte formellement de
Morphisms of Spaces, Lemmas
\ref{spaces-morphisms-lemma-properties-normalization} et
\ref{spaces-morphisms-lemma-normalization-in-integral}.
\end{proof}
